\section{Rogers' proof of the quadratic transformations}\label{rpss}

Since Rogers' papers are not easy to read, we will provide a slightly modified
version of his proof of the quadratic transformations \eqref{rtf}.
Although the notation was not used by Rogers, we find it instructive to introduce the $q$-exponential functions
\begin{align*}
e_q(x)&=\frac 1{(x;q)_\infty}=\sum_{k=0}^\infty\frac{x^k}{(q;q)_k},\qquad
E_q(x)=(-x;q)_\infty=\sum_{k=0}^\infty\frac{q^{k(k-1)/2}x^k}{(q;q)_k}.
\end{align*}
Roger defines the $q$-derivative as
\begin{equation*}(\delta_x f)(x)=\frac{f(x)-f(qx)}{x}. \end{equation*}
It is easy to see that $\delta_x e_q(tx)=t e_q(tx)$,
 and hence $f(\delta_x)e_q(tx)=f(t)e_q(tx)$ for any formal power series $f$. In particular,
\begin{equation}\label{dxe}
e_q(y\delta_x)e_q(tx)=e_q(tx)e_q(ty).\end{equation}
Identifying the coefficient of $t^k$ in \eqref{dxe} gives
\begin{equation}\label{dxh}e_q(y\delta_x)\frac{x^k}{(q;q)_k}=H_k(x,y),\end{equation}
where $H_k$ are defined by the generating function
\begin{equation}\label{eeh}e_q(tx)e_q(ty)= \frac 1{(tx,ty;q)_\infty}=\sum_{k=0}^\infty H_k(x,y)t^k.\end{equation}
In \cite{r1}, Rogers proves \eqref{dxh} in a more direct way and uses it to derive \eqref{dxe}.

The second $q$-exponential function satisfies $\delta_x E_q(tx)=t E_q(qtx)$,
which gives $\delta_x^k E_q(tx)=q^{k(k-1)/2}t^kE_q\bigl(q^ktx\bigr)$ and
\begin{equation*} e_q(y\delta_x)E_q(tx)=\sum_{k=0}^\infty\frac{q^{k(k-1)/2}y^kt^k}{(q;q)_k}\,E_q(q^ktx)
=(-tx;q)_\infty\sum_{k=0}^\infty\frac{q^{k(k-1)/2}y^kt^k}{(-tx,q;q)_k}.\end{equation*}
On the other hand, \eqref{dxh} gives
\begin{equation*}e_q(y\delta_x)E_q(tx)=\sum_{k=0}^\infty t^kq^{k(k-1)/2} H_k(x,y). \end{equation*}
Specializing $t=q$, we obtain \cite[Section~6, equation~(1)]{r2} in the form
\begin{equation}\label{rhi}\sum_{k=0}^\infty q^{k(k+1)/2} H_k(x,y)=(-qx;q)_\infty\sum_{k=0}^\infty\frac{q^{k(k+1)/2}y^k}{(-qx,q;q)_k}.
\end{equation}
By symmetry, one can interchange $x$ and $y$ on the right-hand side.


Roger now observes that the case $y=xq^{1/2}$ of \eqref{eeh} is
\begin{equation*}
 \sum_{k=0}^\infty H_k\bigl(x,xq^{1/2}\bigr)t^k=\frac 1{\bigl(tx;q^{1/2}\bigr)_\infty},\end{equation*}
which implies
\begin{equation*}
H_k\bigl(x,xq^{1/2}\bigr)=\frac{ x^k}{\bigl(q^{1/2};q^{1/2}\bigr)_k}.\end{equation*}
Using this in \eqref{rhi} and replacing $q$ by $q^2$ gives
\begin{subequations}\label{qt}
\begin{gather}
\sum_{k=0}^\infty \frac {q^{k(k+1)}x^k}{(q;q)_k}=\bigl(-q^2x;q^2\bigr)_\infty\sum_{k=0}^\infty\frac{q^{k(k+2)}x^k}{\bigl(-q^{2}x,q^2;q^2\bigr)_k}.
\end{gather}
Interchanging the role of $x$ and $y$ gives the alternative expression
\begin{gather}
\sum_{k=0}^\infty \frac {q^{k(k+1)}x^k}{(q;q)_k}=\bigl(-q^3x;q^2\bigr)_\infty\sum_{k=0}^\infty\frac{q^{k(k+1)}x^k}{\bigl(-q^{3}x,q^2;q^2\bigr)_k}.
\end{gather}
Finally, the case $y=-x$ of \eqref{eeh} is
\begin{equation*}
\sum_{k=0}^\infty H_k(x,-x)t^k=\frac 1{\bigl(t^2x^2;q^{2}\bigr)_\infty},
\end{equation*}
which gives
\begin{equation*}
H_{2k}(x,-x)=\frac{x^{2k}}{\bigl(q^2;q^2\bigr)_k},\qquad H_{2k+1}(x,x)=0.
 \end{equation*}
Using this in \eqref{rhi} gives
\begin{equation}
\sum_{k=0}^\infty \frac{q^{2k(2k+1)}x^{2k}}{\bigl(q^2;q^2\bigr)_k}=(-qx;q)_\infty\sum_{k=0}^\infty\frac{(-1)^kq^{k(k+1)/2}x^k}{(-qx,q;q)_k}.
\end{equation}
\end{subequations}
The identities \eqref{qt} are indeed equivalent to \eqref{rtf}.
